%!TEX root = ../main.tex
% ANEXO I — Repositorios de código y material externo

\chapter{Repositorios de código fuente}\label{anx:repositorios}

\begin{guia}
Los anexos contienen material que no fue producido por el autor o que
es externo al documento: repositorios, datasets de terceros, normativa,
manuales del proveedor. Para cada repositorio: enlace, tecnología,
contenido y estado (activo, archivado).
\end{guia}

\section{GraphSAST}

\prettygitlink{MaxiZk/graphsast}

Monorepo en TypeScript con dos paquetes: \texttt{core}, que contiene el
motor de análisis, la línea de comandos y la API local, y \texttt{viz},
que contiene la interfaz web. Incluye el catálogo de detección en
\texttt{catalog/}, el banco sintético, el evaluador de advisories y los
materiales de la validación con usuarios en
\texttt{docs/validacion/h3/}. La instalación y la ejecución se describen
en el archivo \texttt{README.md}. La versión publicada se encuentra en
\url{https://graph-sast.vercel.app/}.

\completar{Licencia (riesgo R-10) y etiqueta de la versión entregada.}
Estado: activo.

\section{OWASP NodeGoat}

\prettygitlink{OWASP/NodeGoat}

Aplicación Node.js deliberadamente vulnerable, construida por la
fundación OWASP con fines didácticos y distribuida con licencia
Apache-2.0. Se utiliza como corpus de validación sobre código real en el
capítulo~\ref{cap:resultados}, en el commit
\texttt{c5cb68a} del 21 de junio de 2023.
